\section{Final Considerations}
\label{sec:conclusion}

This experience report described the design and the qualitative evaluation of a web-based educational game that supports the understanding of Nielsen's ten heuristics through paired interactive scenarios. The contribution is threefold: the artefact itself, a classroom-applicable procedure for using it in an HCI evaluation course, and a qualitative reading of how a sample of thirty participants described the support the game provided. Participants identified the paired bad/good contrast and the interactive nature of the scenarios as the dominant mechanisms of support, with the reveal screen labelling the concept that links experience to definition. The framing the participants used --- active versus passive, error versus correct, with versus without --- mirrors the experiential and active-learning ideas the artefact was designed around.

The design choices that shape the artefact also shape its limits. With ten heuristics to cover, the absence of a scored failure state and of explicit recognition steps was a deliberate trade-off to keep the total session short enough for classroom use; the bad scenarios still produce real in-task friction, but the player is never asked to commit to identifying the heuristic. The data suggests that this trade was right for the audience studied. The same choices produced the most consistent critique in the data, however: several participants read the artefact as a guided demonstration rather than as a game proper, and asked for the very elements that the design had deferred. The artefact is best understood as a content-delivery and reflection-scaffolding device that complements an instructor, not as a stand-alone teaching system or as an assessment tool.

A short list of next steps follows directly from the participants' feedback. The most cited proposal is to add a recognition-checking mode in which the player attempts to identify the heuristic before the reveal, with stars or lives serving as a real challenge layer. Narration ergonomics --- an in-game mute toggle and shorter, opt-in speech segments --- are inexpensive and would resolve the consistent complaints reported in Section~\ref{sec:results:t4}. A second scenario per heuristic would give players a way to consolidate the contrast, as one participant directly suggested. Beyond the artefact, three larger evaluation steps are open: a delayed re-evaluation that measures retention, a multi-site classroom deployment that broadens the participant profile beyond Computer Science and Software Engineering, and a complementary round in which HCI educators play the game themselves to evaluate it as potential adopters --- a perspective the present study did not capture, since the instructor in the classroom deployment applied the game rather than playing it. We hope these refinements help the artefact continue to support, rather than replace, the work an HCI instructor does in introducing the heuristics to learners.
